% Part: counterfactuals
% Chapter: introduction
% Section: material-conditional

\documentclass[../../../include/open-logic-section]{subfiles}

\begin{document}

\olfileid{cnt}{int}{mat}

\olsection{वास्तविक अभिव्यंजन}

इंग्रजीत सशर्त विधानाचे सर्वात सोपे रूप
‘‘जर \dots तर \dots’’ असे असते; येथे
\dots{} यांच्या जागी स्वतः विधाने येतात,
उदाहरणार्थ, ‘‘जर हे काम सेवकाने केले
असेल, तर माळी निर्दोष आहे.’’ प्राथमिक
तर्कशास्त्रात अशी सशर्त विधाने
$\lif$ या कारकाने दर्शवायला शिकतो:
\dots{} ने दाखवलेले भाग, उदाहरणार्थ,
!!{formula}s $!A$ आणि~$!B$ यांनी
दर्शवायचे; मग संपूर्ण सशर्त विधान
$!A \lif !B$ असे दर्शवतात.

$\lif$ हा कारक \emph{सत्यता-फलनात्मक}
आहे: $!A \lif !B$ चे सत्यतामूल्य---$\True$
किंवा $\False$---हे $!A$ आणि~$!B$
यांच्या सत्यतामूल्यांवरून ठरते.
$!A$ असत्य असेल किंवा $!B$ सत्य
असेल तेव्हा आणि केवळ तेव्हाच $!A \lif !B$
सत्य असते; अन्यथा असत्य असते.
$\pAssign v$ या सत्यतामूल्य-नियुक्तीच्या
संदर्भात $\pSat{v}{!A \lif !B}$
तेव्हा आणि केवळ तेव्हाच, जेव्हा
$\pSat/{v}{!A}$ किंवा
$\pSat{v}{!B}$, अशी व्याख्या
करतो. ही चिन्हार्थमीमांसा
असलेल्या $\lif$ कारकाला
\emph{वास्तविक अभिव्यंजन}
म्हणतात.

या व्याख्येवरून काही प्राथमिक तार्किक
तथ्ये मिळतात. सर्वप्रथम,
वास्तविक अभिव्यंजनासाठी निगमन प्रमेय लागू होते:
\begin{align}
  \text{जर } \Gamma, !A \Entails !B \text{ तर } \Gamma \Entails !A \lif !B
\end{align}
ते सत्यता-फलनात्मक आहे:
$!A \lif !B$ आणि
$\lnot !A \lor !B$ सममूल्य आहेत:
\begin{align}
  !A \lif !B & \Entails \lnot !A \lor !B\\
  \lnot !A \lor !B & \Entails !A \lif !B
  \intertext{वास्तविक अभिव्यंजन त्याच्या
    उत्तरांगावरून आणि पूर्वांगाच्या
    निषेधावरून तार्किकरीत्या निष्पन्न होते:}
  !B & \Entails !A \lif !B\\
  \lnot !A & \Entails !A \lif !B
  \intertext{असत्य वास्तविक अभिव्यंजन
    हे त्याच्या पूर्वांगाचा आणि उत्तरांगाच्या
    निषेधाचा संधियोग याच्याशी सममूल्य असते:
    $!A \lif !B$ असत्य असेल, तर
    $!A \land \lnot !B$ सत्य असते,
    आणि उलटही:}
  \lnot(!A \lif !B) & \Entails !A \land \lnot !B\\
  !A \land \lnot !B & \Entails \lnot(!A \lif !B)
  \intertext{वास्तविक अभिव्यंजनामुळे
    विध्यात्मक अनुमान लागू होते:}
  !A, !A \lif !B & \Entails !B
  \intertext{वास्तविक अभिव्यंजन
    निष्कर्षांचे संयुग्मीकरण करते:}
  !A \lif !B,  !A \lif !C & \Entails !A \lif (!B \land !C)
  \intertext{पूर्वांग अधिक बलवान
    नेहमी करता येते; म्हणजे
    हे अभिव्यंजन \emph{एकस्वनिक} आहे:}
  !A \lif !B & \Entails (!A \land !C) \lif !B
  \intertext{वास्तविक अभिव्यंजन
    संक्रमक आहे; म्हणजे साखळी-नियम वैध आहे:}
  !A \lif !B, !B \lif !C & \Entails !A \lif !C
  \intertext{वास्तविक अभिव्यंजन
    त्याच्या प्रतिपरिवर्तित विधानाशी सममूल्य आहे:}
  !A \lif !B & \Entails \lnot !B \lif \lnot !A\\
  \lnot !B \lif \lnot !A & \Entails !A \lif !B
\end{align}

हे सर्व निष्कर्ष गणिती युक्तिवादात
उपयुक्त आणि निर्विवाद आहेत.
मात्र गणिताबाहेरील संदर्भांतील
सममूल्य निष्कर्षांना कथित
प्रतिदृष्टांत तत्त्वज्ञान आणि भाषाविज्ञान
यांच्या साहित्यात विपुल आहेत.
त्यांवरून असे सुचते की
इंग्रजीतील ‘‘जर \dots तर \dots’’
विधाने दर्शवताना
वास्तविक अभिव्यंजन~$\lif$
हा कारक योग्य नसतो---किंवा
किमान नेहमीच योग्य नसतो.

\end{document}
